\section{Contact resources and their equality cases}
\label{app:primary-contact-resource}

Actual first-jet contacts give a nonnegative resource for each
interpolant factor, in every characteristic. Its zero case is exactly
a power of a received-word graph. These statements concern factors;
relating their charges to normal-divisor costs remains a separate step.

Let $K$ be any field, let $x_1,\ldots,x_n\in K$ be distinct, and let
$r_i\in K$ be received symbols. Fix $w\ge2$, give $(X,Y,R)$ weights
$(1,w,w-1)$, and write $\operatorname{wt}$ for weighted degree.
For $0\ne F\in K[X,Y,R]$, define
\[
 a_i(F)=\operatorname{ord}_t F(x_i+t,r_i+tR+t^2E,R),\qquad
 \epsilon(F)=\operatorname{wt}(F)-\frac wn\sum_i a_i(F).
\]
The substitution is injective after inverting $t$, so contacts are
finite and additive on products. Weighted degrees, and hence charges,
are additive as well. The coefficient field may be $k(Z)$.

\begin{lemma}[Subset contact bound]
\label{lem:contact-subset-average}
If $w\le n$, $W=\operatorname{wt}(F)$, and $0\le u_i\le a_i(F)$ are
integers, then
\begin{equation}
 \#\left\{T\subseteq[n]:|T|=w,\ \sum_{i\in T}u_i>W\right\}
 \le\left\lfloor\frac W{w-1}\right\rfloor.
 \label{eq:contact-subset-bound}
\end{equation}
\end{lemma}
\begin{proof}
For each qualifying $T$, interpolate the symbols by $P_T$ of degree
less than $w$, and put $H_T=\prod_{i\in T}(X-x_i)$.
Contact makes $F(X,P_T+cH_T,P_T'+cH_T')$ divisible by
$\prod_{i\in T}(X-x_i)^{u_i}$. Its degree in $X$ is at most $W$, so it
vanishes identically, with $c$ indeterminate. Consequently
\[
 R-P_T'-\frac{H_T'}{H_T}(Y-P_T)
\]
is a factor over $K(X)$. Distinct $T$ give distinct factors: their
slopes have simple poles of residue one at precisely their respective
nodes, in every characteristic. Their number is at most the total
jet degree, which is at most $\lfloor W/(w-1)\rfloor$.
\end{proof}

\begin{theorem}[Universal contact resource]
\label{thm:contact-universal-resource}
Suppose $n\ge2w$. Write $r=\deg_R F$, $A_r=[R^r]F$,
$q=\deg_Y A_r$, and $B=[Y^q]A_r\ne0$, with $d=\deg_X B$.
Then, in every characteristic,
\begin{equation}
 a_i(F)\le q+2\operatorname{ord}_{x_i}B,\qquad
 \epsilon(F)\ge r(w-1)+\left(1-\frac{2w}{n}\right)d
 \ge r(w-1)\ge0.
 \label{eq:contact-average-bound}
\end{equation}
\end{theorem}
\begin{proof}
Fix a node and translate $Y$ by its received symbol. This preserves
$r,q,B$. In the translated polynomial take the homogeneous
$(Y,R)$-degree $q+r$ piece $G$, whose $Y^qR^r$ coefficient is
$B(x_i+t)$. The $R$-degree bound forces $G=Y^qH$, where $H$ is
homogeneous of degree $r$ and $H(t,0,1)=B(x_i+t)$.
Distinct homogeneous pieces remain distinct homogeneous pieces in
$(R,E)$ after contact substitution, so they cannot cancel. Thus
\[
 a_i(F)\le\operatorname{contact}(G)
 =q+\operatorname{contact}(H).
\]
Expand $h(t,Y)=H(t,Y,1)=\sum_j c_j(t)(Y-t)^j$.
Homogeneity and the substitution $Y=tR+t^2E$ give
\[
 b:=\operatorname{contact}(H)=\min_j\{\operatorname{ord}_t c_j+2j\}.
\]
But $B(x_i+t)=h(t,0)=\sum_j c_j(t)(-t)^j$.
Each summand has order at least $\lceil b/2\rceil$, since
$\operatorname{ord}_t c_j\ge0$. Hence $b\le2\operatorname{ord}_{x_i}B$.
Summing over the distinct nodes yields $\sum_i a_i(F)\le nq+2d$.
Finally $\operatorname{wt}(F)\ge r(w-1)+qw+d$, proving the claim.
\end{proof}

\begin{corollary}[Additive primary and uniform bounds]
\label{cor:contact-factor-resource}
\label{thm:contact-regular-gap}
Assume $n\ge2w$. If $Q_0=c\prod_j F_j^{e_j}\ne0$ has contact at
least $m$ at every node and weight $W_0<mA$, then
\begin{equation}
 (w-1)\sum_j e_j\deg_R F_j
 \le\sum_j e_j\epsilon(F_j)<m(A-w).
 \label{eq:contact-additive-resource}
\end{equation}
If instead $F\ne0$ has weight at most $aw+e$, where $a,e\ge0$ are
integers, and contact at least $a$ at every node, then
\[
 \deg_R F\le\left\lfloor\frac e{w-1}\right\rfloor.
\]
In particular $e<w-1$ forces complete independence of $R$, including
in small characteristic.
\end{corollary}
\begin{proof}
The theorem and additivity give the first lower bound, while
\[
 \sum_j e_j\epsilon(F_j)=W_0-\frac wn\sum_i a_i(Q_0)
 \le W_0-mw<m(A-w).
\]
For the uniform statement, $(w-1)\deg_R F\le\epsilon(F)\le e$.
\end{proof}

\begin{theorem}[Zero-charge rigidity in every characteristic]
\label{cor:contact-zero-charge}
\label{thm:contact-uniform-rigidity}
Suppose $n\ge2w$. Then $\epsilon(F)=0$ if and only if
\[
 F=c(Y-P(X))^q,\qquad c\in K^\times,
 \quad\deg P\le w,\quad P(x_i)=r_i\ \text{for all }i,
\]
for an integer $q\ge1$, or $F$ is a nonzero scalar ($q=0$).
Consequently, weight at most $aw$ and contact at least $a$ at every
node force this form with $q=a$.
\end{theorem}
\begin{proof}
The universal bound first gives $\deg_R F=0$. Put $q=\deg_Y F$ and
let $B(X)$ be its leading coefficient, of degree $d$.
The coefficient of $R^qE^0$ in the contact substitution gives the
stronger bound $a_i(F)\le q+\operatorname{ord}_{x_i}B$. Therefore
\[
 \epsilon(F)\ge\operatorname{wt}(F)-qw-\frac wn d
 \ge\left(1-\frac wn\right)d.
\]
Equality forces $d=0$, $\operatorname{wt}(F)=qw$, and $a_i(F)=q$ at
every node. The case $q=0$ is scalar.

First suppose $q$ is nonzero in $K$. Write
$F=BY^q+C(X)Y^{q-1}+\cdots$, with $B\in K^\times$ and $\deg C\le w$.
The coefficient of $t^{q-1}R^{q-1}$ gives $qBr_i+C(x_i)=0$.
Thus $P=-C/(qB)$ has degree at most $w$ and matches all symbols.
Translate $Y$ by $P(X)$. Weight is preserved and the new polynomial
has zero-word contact $q$: locally, the difference
$P(x_i+t)-P(x_i)$ is absorbed by an invertible translation of $R$.
If $C_j(X)$ is its $Y^j$ coefficient for $j<q$, contact implies
\[
 \prod_i(X-x_i)^{q-j}\mid C_j,
 \qquad\deg C_j\le(q-j)w<(q-j)n.
\]
Hence every such coefficient vanishes, proving the required form.

It remains to treat characteristic $p\mid q$. Differentiating the
contact identity in $R$ shows that $F_Y$ has contact at least $q-1$,
since the derivative is $t(F_Y)_{\rm sub}$. If $F_Y\ne0$, its weight
is at most $(q-1)w$. Nonnegative charge forces equality and zero
charge. The preceding characteristic-free argument for $R$-free
zero-charge polynomials would then give
$\deg_Y F_Y=q-1$, whereas the leading term differentiates to zero
and $\deg_Y F_Y\le q-2$. Thus $F_Y=0$.

Differentiating in $t$ now gives contact at least $q-1$ for $F_X$.
If $F_X\ne0$, let $j$ be its $Y$-degree and $B_j(X)$ its leading
coefficient. All $Y$-exponents of $F$ are divisible by $p$, and its
leading coefficient is constant, so $j\le q-p$. Put $d'=q-j\ge2$.
The $R$-free contact bound forces
\[
 \deg B_j\ge n(d'-1),\qquad
 \deg B_j\le d'w-1.
\]
This is impossible, since $n\ge2w$ implies $n(d'-1)\ge d'w$.
Therefore $F_X=0$ too. Over the perfect closure of $K$, write $F=G^p$.
Then $G$ has weight $(q/p)w$ and contact $q/p$ at the same nodes.
Repeat until its $Y$-degree is nonzero in the field, apply the previous
case, and raise back to the appropriate power. The resulting $P$
descends to $K$: interpolation at any $w+1$ original nodes uniquely
recovers it from the original symbols.

Conversely, the displayed graph power has weight $qw$ and exact
contact $q$. For the final assertion, nonnegative charge and the two
bounds force zero charge, weight $aw$, and contact $a$.
\end{proof}

For nodes in $k$ and affine symbols $r_i=f_i+Zg_i$, interpolation at
$w+1$ nodes gives $P=P_0+ZP_1$, with $\deg P_0,\deg P_1\le w$.
Thus over $k(Z)$ every positive-weight irreducible zero-charge factor
is a received-line graph factor. If the original polynomial lies in
$k[Z,X,Y,R]$, the scalar in its graph-power expression is its leading
$Y$-coefficient and lies in $k[Z]$.

\paragraph{Finite scope and the remaining geometric step.}
At $n=262144$, $w=131071$, $m=118$, and $A=181275$, the total charge
is strictly below $m(A-w)=5924072$. The resulting total $R$-degree
bound is $45$, weaker than the existing source cap $36$; it gives no
benchmark improvement. Nor does the resource alone exclude the
binding carrier flags $(\deg_R,\deg_{Y,R},\deg_{Y,R,Z})=(12,55,3261)$.
Writing $T=Y-Z$ and $\Lambda=\prod_i(X-x_i)$, the polynomial
\[
 F=T^{55}+\Lambda^{10}(R^{12}+R+Z^{3261})
       -\Lambda^9\Lambda'T
\]
has these flags, weight $55w$, and exact contact ten for the word $Z$.
Its charge is $45w=5898195$, within the budget, and
$Q_0=FT^{108}$ has contact $118$ and weight $163w<118A$.
This is source compatibility, not a bad-label construction.
A first-tail component is a codimension-two intersection component,
not another primary factor: charging each such component separately
requires a new inequality controlling their aggregate normal-divisor
cost from the carrier's single charge.
